% Options for packages loaded elsewhere
\PassOptionsToPackage{unicode}{hyperref}
\PassOptionsToPackage{hyphens}{url}
%
\documentclass[
  11pt,
]{article}
\usepackage{amsmath,amssymb}
\usepackage{iftex}
\ifPDFTeX
  \usepackage[T1]{fontenc}
  \usepackage[utf8]{inputenc}
  \usepackage{textcomp} % provide euro and other symbols
\else % if luatex or xetex
  \usepackage{unicode-math} % this also loads fontspec
  \defaultfontfeatures{Scale=MatchLowercase}
  \defaultfontfeatures[\rmfamily]{Ligatures=TeX,Scale=1}
\fi
\usepackage{lmodern}
\ifPDFTeX\else
  % xetex/luatex font selection
\fi
% Use upquote if available, for straight quotes in verbatim environments
\IfFileExists{upquote.sty}{\usepackage{upquote}}{}
\IfFileExists{microtype.sty}{% use microtype if available
  \usepackage[]{microtype}
  \UseMicrotypeSet[protrusion]{basicmath} % disable protrusion for tt fonts
}{}
\makeatletter
\@ifundefined{KOMAClassName}{% if non-KOMA class
  \IfFileExists{parskip.sty}{%
    \usepackage{parskip}
  }{% else
    \setlength{\parindent}{0pt}
    \setlength{\parskip}{6pt plus 2pt minus 1pt}}
}{% if KOMA class
  \KOMAoptions{parskip=half}}
\makeatother
\usepackage{xcolor}
\usepackage[margin=25mm]{geometry}
\usepackage{color}
\usepackage{fancyvrb}
\newcommand{\VerbBar}{|}
\newcommand{\VERB}{\Verb[commandchars=\\\{\}]}
\DefineVerbatimEnvironment{Highlighting}{Verbatim}{commandchars=\\\{\}}
% Add ',fontsize=\small' for more characters per line
\newenvironment{Shaded}{}{}
\newcommand{\AlertTok}[1]{\textcolor[rgb]{1.00,0.00,0.00}{\textbf{#1}}}
\newcommand{\AnnotationTok}[1]{\textcolor[rgb]{0.38,0.63,0.69}{\textbf{\textit{#1}}}}
\newcommand{\AttributeTok}[1]{\textcolor[rgb]{0.49,0.56,0.16}{#1}}
\newcommand{\BaseNTok}[1]{\textcolor[rgb]{0.25,0.63,0.44}{#1}}
\newcommand{\BuiltInTok}[1]{\textcolor[rgb]{0.00,0.50,0.00}{#1}}
\newcommand{\CharTok}[1]{\textcolor[rgb]{0.25,0.44,0.63}{#1}}
\newcommand{\CommentTok}[1]{\textcolor[rgb]{0.38,0.63,0.69}{\textit{#1}}}
\newcommand{\CommentVarTok}[1]{\textcolor[rgb]{0.38,0.63,0.69}{\textbf{\textit{#1}}}}
\newcommand{\ConstantTok}[1]{\textcolor[rgb]{0.53,0.00,0.00}{#1}}
\newcommand{\ControlFlowTok}[1]{\textcolor[rgb]{0.00,0.44,0.13}{\textbf{#1}}}
\newcommand{\DataTypeTok}[1]{\textcolor[rgb]{0.56,0.13,0.00}{#1}}
\newcommand{\DecValTok}[1]{\textcolor[rgb]{0.25,0.63,0.44}{#1}}
\newcommand{\DocumentationTok}[1]{\textcolor[rgb]{0.73,0.13,0.13}{\textit{#1}}}
\newcommand{\ErrorTok}[1]{\textcolor[rgb]{1.00,0.00,0.00}{\textbf{#1}}}
\newcommand{\ExtensionTok}[1]{#1}
\newcommand{\FloatTok}[1]{\textcolor[rgb]{0.25,0.63,0.44}{#1}}
\newcommand{\FunctionTok}[1]{\textcolor[rgb]{0.02,0.16,0.49}{#1}}
\newcommand{\ImportTok}[1]{\textcolor[rgb]{0.00,0.50,0.00}{\textbf{#1}}}
\newcommand{\InformationTok}[1]{\textcolor[rgb]{0.38,0.63,0.69}{\textbf{\textit{#1}}}}
\newcommand{\KeywordTok}[1]{\textcolor[rgb]{0.00,0.44,0.13}{\textbf{#1}}}
\newcommand{\NormalTok}[1]{#1}
\newcommand{\OperatorTok}[1]{\textcolor[rgb]{0.40,0.40,0.40}{#1}}
\newcommand{\OtherTok}[1]{\textcolor[rgb]{0.00,0.44,0.13}{#1}}
\newcommand{\PreprocessorTok}[1]{\textcolor[rgb]{0.74,0.48,0.00}{#1}}
\newcommand{\RegionMarkerTok}[1]{#1}
\newcommand{\SpecialCharTok}[1]{\textcolor[rgb]{0.25,0.44,0.63}{#1}}
\newcommand{\SpecialStringTok}[1]{\textcolor[rgb]{0.73,0.40,0.53}{#1}}
\newcommand{\StringTok}[1]{\textcolor[rgb]{0.25,0.44,0.63}{#1}}
\newcommand{\VariableTok}[1]{\textcolor[rgb]{0.10,0.09,0.49}{#1}}
\newcommand{\VerbatimStringTok}[1]{\textcolor[rgb]{0.25,0.44,0.63}{#1}}
\newcommand{\WarningTok}[1]{\textcolor[rgb]{0.38,0.63,0.69}{\textbf{\textit{#1}}}}
\setlength{\emergencystretch}{3em} % prevent overfull lines
\providecommand{\tightlist}{%
  \setlength{\itemsep}{0pt}\setlength{\parskip}{0pt}}
\setcounter{secnumdepth}{-\maxdimen} % remove section numbering
\ifLuaTeX
  \usepackage{selnolig}  % disable illegal ligatures
\fi
\usepackage{bookmark}
\IfFileExists{xurl.sty}{\usepackage{xurl}}{} % add URL line breaks if available
\urlstyle{same}
\hypersetup{
  pdftitle={P16.4: positive graph interpolation and the remaining actual curve bridge},
  pdfauthor={Unverified research continuation; relative to six complete caller grants},
  hidelinks,
  pdfcreator={LaTeX via pandoc}}

\title{P16.4: positive graph interpolation and the remaining actual
curve bridge}
\author{Unverified research continuation; relative to six complete
caller grants}
\date{7 October 2026}

\begin{document}
\maketitle

\section{Status and exact unchanged
domain}\label{status-and-exact-unchanged-domain}

This manuscript does \textbf{not} close original P16.4. Its new
unconditional mathematical result is the finite graph interpolation
theorem below, with a complete proof. It also proves a conditional
actual curve equation, including loops, arbitrary Radon support,
all-point uniqueness, the raw constants, and finite-orbit descent. The
geometric antecedents of that conditional theorem have not all been
verified under the requested recursive source rule. In particular, this
submission does not assert an unconditional all-genus equation. The
original arbitrary-dimensional equation remains open in this submission.
Compilation, source retrieval, and finite arithmetic checks are not
mathematical certification.

The unchanged original domain is a smooth geometrically integral
projective variety \(X\) of arbitrary dimension over a complete
characteristic-zero real-valued non-Archimedean field \(K\) with residue
characteristic zero, and an actual ample rational line bundle \(L\).
Trivial, nontrivial, discrete, nondiscrete, and non-algebraically-closed
fields are included. Write \(T=X^{\mathrm{an}}\) and \(V=(L^n)>0\).
Targets have raw mass \(V\). All metrics are actual all-point decreasing
Fubini--Study net endpoints on the actual bundle, including literal pole
values. P16.3, P16.5, and P16.7 are separate axes.

For a continuous psh reference \(\psi\), scalars mean absolute metrics
\(\psi+u\). The continuous energy and its extension are \[
 E_\psi(u)=\frac1{n+1}\sum_{j=0}^n\int_Tu\,\operatorname{MA}(u^{[j]},0^{[n-j]}),\qquad
 E_\psi(u)=\inf_{v\geq u,\ v\text{ continuous psh}}E_\psi(v).
\] The first formula initially applies to continuous potentials; the
second defines extended energy on actual singular endpoints. The class
\(\mathcal E^1\) consists of endpoints of finite extended energy. The
target dual is \[
 J_\psi(\mu)=\sup_{v\text{ continuous psh}}\left(E_\psi(v)-\int_Tv\,d\mu\right),\qquad \mu(T)=V.
\] Constants cancel, since \(E(u+c)=E(u)+Vc\). An actual reference
change adds a finite constant to this dual. No measure is defined by
naming an inverse target.

\section{The six complete caller
prerequisites}\label{the-six-complete-caller-prerequisites}

Only the complete supplied bodies O217, D214, N220, T171, Q222, and
C218, including every supplied harmless clarification, are mathematical
grants. Their historical draft headings do not override the caller's
present accepted status. These labels are not registered supporting note
IDs. This submission's emitted note n1 indexes their permitted use;
subsequent supporting IDs refer only to emitted notes.

O gives the genuine actual mixed \(\mathcal E^1\) operator over the
entire original field and geometric domain, raw mass \(V\),
integrability of every actual \(\mathcal E^1\) test against every mixed
output, arbitrary directed weighted decreasing continuity, extended
decreasing-energy continuity including \(-\infty\), reference
covariance, affine curvature calculus on admissible convex combinations,
the segment polynomial and derivative, and complete-field pushforward.
We use extended energy before asserting endpoint membership. Zero
coefficients at poles are omitted. The supplied normalization correction
handles nonpositive endpoints by shifting the reference and then letting
the shift tend to zero; it is not a constrained regularization
assumption.

D gives the jointly continuous completed-closure Galois action, the
all-point quotient, actual invariant psh descent, extended energy
equality, the unique invariant Radon lift, continuous-dual equality, and
the continuous field laws. Its equation implication requires actual
upstairs existence and uniqueness; it is not an existence theorem. Raw
mass is preserved, without any field-degree factor.

N gives all-\(\mathcal E^1\) integrability if and only if finite
continuous dual \textbf{only after} a continuous actual reference with
\(\operatorname{MA}(0)=V\delta_\xi\) and \(\sup u=u(\xi)\) for every FS
scalar has been established. It cannot be applied to an unverified
abstract graph reference. It gives coercivity, not attainment.

T gives the entire actual ample equation and every-solution all-point
uniqueness over every trivially valued characteristic-zero field. Q
gives the split projective-line equation, its actual disk atom metrics,
its local disk classification, and every-solution all-point uniqueness.
Q supplies no arbitrary-genus skeleton or model geometry. C gives strict
Borel locality, open locality, and mixed and pure comparison for the
genuine operator in O's domain. Comparison alone is not equation
existence or every-solution uniqueness.

The complete supplied bodies are preserved as
\texttt{sources/grant-O.tex}, \texttt{grant-D.tex},
\texttt{grant-N.tex}, \texttt{grant-T.tex}, \texttt{grant-Q.tex}, and
\texttt{grant-C.tex}. The unchanged \texttt{input.json} preserves the
complete caller input and clarifications. No historical review output or
another campaign was read.

\section{Positive interpolation on graphs with irrational
cycles}\label{positive-interpolation-on-graphs-with-irrational-cycles}

\textbf{Theorem 1 (finite graph statement).} Let
\(\Lambda\subset\mathbb R\) be a nonzero divisible additive subgroup.
Let \(G\) be a finite connected compact metric graph, with all edge
lengths in \(\Lambda\), and with positive fixed weights \(w_e\). Let
\(\theta\) be a signed vertex measure of total mass \(V>0\), and let
\(\mu\) be a nonnegative vertex measure of the same mass. Fix a vertex
\(z\). There is a unique continuous function \(F\), affine on each edge,
satisfying \[
 F(z)=0,\qquad \theta+\Delta F=\mu,
\] where \(\Delta\) is the sum of weighted outgoing slopes. There are
finite subdivisions and functions \(G_i\) with rational slopes, vertex
values in \(\Lambda\), and one common denominator at each stage, such
that \[
 G_i\longrightarrow F\text{ uniformly},\qquad
 \theta+\Delta G_i\geq0,\qquad
 \theta+\Delta G_i\longrightarrow\mu\text{ weakly}.
\] Every original vertex and every inserted bend has strictly positive
curvature at each stage. This is a graph theorem, not a psh recognition
theorem.

At an oriented interior bend with left slope \(q_1\) and right slope
\(q_2\), the curvature atom is \(w_e(q_2-q_1)\). At the initial endpoint
the contribution is \(w_eq_1\); at the terminal endpoint it is
\(-w_eq_2\). Thus an upward bend is positive and decreases the total
curvature retained at the two old endpoints. All three contributions
have sum zero. This is the positive outgoing-slope convention. KRZB's
graph divisor uses its negative.

Subdivide a self-loop once at length \(\ell/2\in\Lambda\), extending
\(\theta\) and \(\mu\) by zero at the new vertex. This allows a proof
with distinct endpoints for every edge; loops are not deleted. A
one-vertex graph with no edges has \(\theta=\mu=V\delta_z\), and
\(F=G_i=0\) suffices.

For an edge-affine function \(H\), direct cancellation of edge
contributions gives \[
 -\sum_vH(v)\Delta H(v)=\sum_{e=ab}\frac{w_e}{\ell_e}\bigl(H(b)-H(a)\bigr)^2.
\] Polarization proves symmetry of the vertex matrix. Connectivity makes
its kernel exactly the constants. Its image is therefore precisely the
zero-sum vectors. Anchoring at \(z\) gives the claimed unique solution
by finite-dimensional linear algebra. The same argument works with
arbitrary positive real edge weights and signed \(\theta\); no
positivity of the starting model is used in this graph step.

Let there be \(k\) original vertices after loop subdivision, and set
\(\rho=(V/k)\sum_v\delta_v\). For \(0<\epsilon<1\) put \[
 \mu_\epsilon=(1-\epsilon)\mu+\epsilon\rho.
\] Every original vertex has a positive curvature margin. The anchored
solution \(F_\epsilon\) to \(\theta+\Delta F_\epsilon=\mu_\epsilon\)
depends linearly on its right-hand side. Consequently
\(F_\epsilon\to F\) uniformly, with convergence of every original edge
slope.

The subgroup \(\Lambda\) is dense: if \(a\ne0\) belongs to it, then
\(\mathbb Qa\subset\Lambda\) is dense. Approximate the finitely many
shared vertex values \(F_\epsilon(v)\) by \(g_v\in\Lambda\), with
\(g_z=0\) and error at most \(h\). For the oriented edge \(a\to b\) put
\(d=g_b-g_a\) and \(c=d/\ell\). The chord slope differs from the slope
of \(F_\epsilon\) by at most \(2h/\ell\).

Choose rational numbers \(q_1<c<q_2\), both within \(\eta/2\) of \(c\),
and arrange \(2h/\min_e\ell_e<\eta/2\). Strict bracketing is possible
even when \(c\) is rational. Set \[
 t=\frac{d-q_2\ell}{q_1-q_2}.
\] Because \(q_1<c<q_2\), one has \(0<t<\ell\). Because \(\Lambda\) is a
rational vector space, \(t\in\Lambda\). Use slope \(q_1\) from \(a\) to
\(t\) and slope \(q_2\) afterwards. The terminal value is exactly
\(g_b\), and the inserted vertex value \(g_a+q_1t\) belongs to
\(\Lambda\). The same \(g_v\) is used on every incident edge;
compatibility around every cycle is exact. No claim that rational slopes
on fixed irrational edges are dense is needed.

Every outgoing slope at an original vertex differs from the
corresponding slope of \(F_\epsilon\) by less than \(\eta\). Choose \[
 \eta\sum_{e\ni v}w_e<\frac12\mu_\epsilon(v)
\] at every original vertex. Its new curvature is then at least half its
positive margin. The single inserted point of each edge has mass
\(w_e(q_2-q_1)>0\). Total mass remains \(V\), since \(\Delta\) has total
mass zero. The edgewise uniform error against \(F_\epsilon\) is at most
\(h+\eta\ell_e\). The total inserted mass is less than
\(\eta\sum_ew_e\), and the old vertex masses converge to those of
\(\mu\) as first the approximation errors and then \(\epsilon\) tend to
zero. A diagonal sequence of these finite choices proves all the stated
convergences. One common integer clears all rational slopes;
multiplication by that integer preserves membership of every vertex
value in \(\Lambda\). Thus the approximants are integral
\(\Lambda\)-affine after clearing denominators.

There is a genuine obstruction to the simpler fixed-edge strategy. On
the two-edge cycle with lengths \(1\) and \(\sqrt2\), an edge-affine
function with shared difference \(d\) and rational slopes on both edges
must have \(d\in\mathbb Q\cap\sqrt2\mathbb Q=\{0\}\). Such functions are
constant. Refining edges with movable positive bends removes this
obstruction.

For the exact prototype \(\theta=(1,1)\), \(\mu=(2,0)\),
\(\epsilon=1/10\), the anchored difference is \(d=(9/10)(2-\sqrt2)\).
Slopes \((52/100,53/100)\) on the length-one edge and
\((37/100,38/100)\) on the length-\(\sqrt2\) edge give bend positions
\(-127+90\sqrt2\) and \(-180+128\sqrt2\), respectively. The old masses
are \(189/100\) and \(9/100\), and each new bend has mass \(1/100\).
Their sum is exactly \(2\). These statements are proved by the formulas
above and independently checked by exact arithmetic.

\section{Finite-core Dirichlet kernels, including
loops}\label{finite-core-dirichlet-kernels-including-loops}

In this section weights are one, as on a split skeleton over a complete
algebraically closed field. Let \(G\) be any finite connected compact
metric graph and \(z\in G\). Subdivide at points being used. Define
\(j_z(x,y)\) by \[
 \Delta_xj_z(x,y)=\delta_z-\delta_y,\qquad j_z(z,y)=0.
\] The preceding matrix proof gives existence and uniqueness, also on
graphs with cycles. Define \[
 D_G(f,h)=\sum_e\int_e f'h'\,dt=-\int_Gf\,\Delta h.
\] The integration formula follows on each affine segment and by
telescoping at vertices. It gives \[
 D_G(j_z(\cdot,x),j_z(\cdot,y))=j_z(x,y)=j_z(y,x).
\] The minimum principle gives \(j_z\geq0\): a negative minimum is not
at \(z\), is incompatible with the negative sink mass at \(y\), and
otherwise propagates through zero outgoing slopes until it reaches \(z\)
or \(y\). The maximum is at \(y\) (or the kernel is zero when \(y=z\)).

Each edge slope of \(j_z(\cdot,y)\) has absolute value at most one. For
a regular level between zero and its maximum, sum \(\Delta j_z\) over
the vertices below that level and the corresponding partial edge
segments. The crossing slopes are nonnegative, their sum is one, and
each is at most one. Every nonconstant edge meets such a level. Symmetry
therefore yields \[
 |j_z(x,y)-j_z(x,y')|\leq d_G(y,y').
\] This also proves joint continuity on \(G^2\). The one-point/no-edge
kernel is zero.

Suppose \(G'\) is obtained by subdividing \(G\) and attaching finite
trees, each at a single point of \(G\). Let \(r:G'\to G\) collapse these
trees. On an attached rooted tree write \(\lambda(x,y)\) for the length
of the common initial part of the two paths from its attachment point;
put it equal to zero unless both points are in that same tree. Direct
outgoing-slope cancellation, including the attachment vertex, gives \[
 j^{G'}_z(x,y)=j^G_z(rx,ry)+\lambda(x,y).
\] Both sides have the same source, sink, and anchor; uniqueness proves
the formula. Adding several trees, or attaching them in stages, gives
the same identity. Consequently the old kernel is unchanged when one
point is in \(G\), and the retracted kernels increase pointwise as
branches are added. This is not an assumption that a graph with loops is
a tree. The loops stay in \(G\). Weighted branches would use resistance
lengths \(\ell/w\), not ordinary lengths; only unweighted upstairs
branches are used below.

We also need a finite-graph maximum-support lemma. Let \(h\) be a kernel
potential with \(\Delta h=\alpha-\sigma\), where \(\alpha,\sigma\) are
positive Radon measures of the same mass and \(\sigma\) is supported on
a nonempty finite set \(S\). Then \(\max_Gh=\max_Sh\). Indeed \(h\) is
Lipschitz with edge derivatives of bounded variation, by the kernel
estimate. Put \(M=\max_Sh\) and \(b=(h-M)^+\). Edgewise Stieltjes
integration by parts and vertex cancellation give \[
 \sum_e\int_e|b'|^2=-\int_G b\,\Delta h
       =-\int_Gb\,d\alpha\leq0.
\] Here \(b=0\) on \(S\), and \(b'=\mathbf1_{\{h>M\}}h'\) almost
everywhere. Thus \(b\) is constant on the connected graph and vanishes
on \(S\). The claim follows. The integration formula can alternatively
be obtained first with polygonal test functions, then by edgewise
uniform and Sobolev approximation to \(b\); the finite measure
\(\Delta h\) passes through uniform convergence. This is elementary
one-dimensional graph integration, not a CLD smoothness assertion about
a Thuillier lisse function.

\section{Explicit conditional geometric
antecedents}\label{explicit-conditional-geometric-antecedents}

The following requirements are \textbf{additional hypotheses}, not
conclusions about every original curve. They isolate the remaining
bridge precisely.

(R1) Over a complete algebraically closed nontrivially valued field
\(C\), \(X\) is a smooth connected projective curve with its actual
ample rational \(L\), \(V=\deg L\). There is a finite semistable core
\(\Sigma_0\), with type-II vertices and edge lengths in
\(\Lambda=-\log|C^*|\), and an actual projective rational model of \(L\)
with metric \(\psi_0\). Its component-degree measure \(\theta\) is
supported on that core and has raw mass \(V\).

(R2) The finite skeleta containing \(\Sigma_0\) form a directed family
of continuous retractions \(r_G:T\to G\). Retractions compose in the
inclusion order and separate all analytic points. Enlargements add only
trees to an old graph and subdivide its edges. Every type-II and
type-III point belongs to some such graph; the complement of a core is a
union of analytic open balls, with the usual disk classification and
topology. For type-II vertices the graphs and refinements have
value-group lengths. These are actual analytic and model assertions, not
just a compact inverse-system declaration.

(R3) If \(G\) is such a graph with value-group vertices and \(f\) is
rational-slope/value-group-valued PA with \(\theta_G+\Delta f\geq0\),
there is a refined \textbf{actual projective rational model on the same
bundle} whose additive metric is \(\psi_0+f\circ r_G\), whose
special-component degrees are the atoms of \(\theta_G+\Delta f\), and
whose genuine continuous Monge--Ampere measure has those exact
coefficients. Here \(\theta_G\) has the old degrees and zero on
contracted new components. The refined bundle is nef, and its metric is
recognized and uniformly approximated by actual fixed-bundle FS metrics.
Real PA positive-curvature limits then pass by the genuine continuous
uniform law. Theorem 1 proves this last approximation step, but it does
not prove the preceding model assertions.

(R4) On each complementary formal-fibre open ball, a local analytic
frame of a cleared power of \(L\) has constant \(\psi_0\)-norm. Adding a
scalar factoring through \(\Sigma_0\) retains that property. Nested
finite disk ancestors converge to every terminal point, and analytic
section norms obey disk inclusion there. This is the needed
terminal-point geometry. It must be verified on actual charts; it is not
granted by calling an abstract graph potential psh.

For original-field descent of the all-test criterion we further require
(R5): for \(X_C\) there is a nonempty finite \(G_K\)-orbit of actual
type-II points. A model defined over a finite extension with
geometrically split components would supply it, but the finite-extension
existence assertion has not been checked recursively here. No fixed
point of a Galois action on a graph with loops is assumed.

\textbf{Theorem 2 (conditional actual all-genus curve theorem).}
Assuming (R1)--(R4) for the algebraically closed carrier, the genuine O
operator solves every positive Radon target of raw mass \(V\) and finite
continuous dual. Every solution in the actual \(\mathcal E^1\) cone
differs by one finite real constant at every analytic point. The
finite-dual condition equals integrability of every actual finite-energy
potential, and the solution attains the variational supremum. If these
hypotheses hold on the completed closure of an original nontrivial
field, D transfers existence and every-solution uniqueness for
finite-dual targets. With (R5), the downstairs all-\(\mathcal E^1\)
criterion also follows from a genuinely descended single norming atom. T
handles trivial fields. This is an implication with unverified geometric
antecedents, not an unconditional all-genus result.

\section{Actual atom metrics under the geometric
antecedents}\label{actual-atom-metrics-under-the-geometric-antecedents}

Fix a type-II vertex \(z\) of \(\Sigma_0\). Theorem 1 solves and
approximates the anchored real PA function \(F_z\) with \[
 F_z(z)=0,\qquad\theta+\Delta F_z=V\delta_z.
\] Apply (R3) to each rational approximation. Its bundle is the actual
old bundle tensored with the actual vertical rational Cartier divisor,
rather than a numerical replacement. Its component degrees are
nonnegative and its metric is actual continuous psh. Uniform limits and
genuine continuous curvature continuity give the actual reference \[
 \psi_z=\psi_0+F_z\circ r_{\Sigma_0},\qquad\operatorname{MA}(\psi_z)=V\delta_z.
\] On a larger finite graph extend \(F_z\) constantly along new trees.
If \(y\) is a type-II vertex, then \[
 F_y=F_z-Vj_z(\cdot,y),\qquad\theta_G+\Delta F_y=V\delta_y.
\] Theorem 1 and (R3) again give an actual continuous atom metric. Its
scalar relative to \(\psi_z\) is precisely \[
 a_y=-Vj_z(\cdot,y),\qquad a_y(z)=0,\qquad\operatorname{MA}(a_y)=V\delta_y.
\] For any type-III \(y\) in an edge choose type-II points tending to
it, using density of \(\Lambda\). The kernel Lipschitz estimate makes
the actual metrics converge uniformly, and \(\delta_{y_i}\to\delta_y\)
in the analytic topology. Thus the same assertion holds at type III. One
does not apply Theorem 1 directly to a subdivision whose type-III edge
lengths need not lie in \(\Lambda\).

The global function \(j_z(\cdot,y)\) is independent of the chosen finite
graph containing \(y\): the branch-extension formula leaves the kernel
unchanged when one argument is in the old graph. Atom metrics are now
actual operators' outputs. Before (R3) this notation is only graph data.

\section{A genuine norming reference and weighted curve
symmetry}\label{a-genuine-norming-reference-and-weighted-curve-symmetry}

Here and below scalars use \(\psi_z\). For actual curve
\(p,q\in\mathcal E^1\), O's energy identity, affine law, and segment
derivative imply \[
 \int_Tp\,\operatorname{MA}(q)-\int_Tq\,\operatorname{MA}(p)
       =V\bigl(p(z)-q(z)\bigr). \tag{1}
\] To verify the constant, expand
\(E((1-t)p+tq)=\frac12\int((1-t)p+tq)\operatorname{MA}((1-t)p+tq)+\frac V2((1-t)p(z)+tq(z))\)
and compare its derivative at zero with O's
\(\int(q-p)\operatorname{MA}(p)\). Each integral is finite by O, and
\(p(z),q(z)\) are finite because O integrates each test against the
reference output \(V\delta_z\). No common pole values are subtracted.

For a continuous psh \(p\) normalize \(p(z)=0\), and use \(q=a_y\).
Equation (1) gives \[
 p(y)=-\int_Tj_z(x,y)\,d\operatorname{MA}(p)(x)\leq0. \tag{2}
\] Finite graph points are dense in \(T\). A direct proof uses the real
algebra generated by continuous functions pulled back from finite
retractions: it separates points by (R2), and hence is uniformly dense
in \(C(T)\) by Stone--Weierstrass. Any fixed element is unchanged by
every sufficiently large retraction. Therefore \(r_Gx\to x\) against
every continuous test, for each \(x\). In particular the union of the
finite graphs is dense. Continuity of \(p\) proves \(\sup_Tp=p(z)\).

Every FS scalar is continuous psh, so this establishes the exact N
hypothesis for \(\psi_z\), on the actual bundle. For an actual
decreasing FS endpoint \(p_i\downarrow p\), the finite numbers
\(p_i(z)=\sup p_i\) decrease. A limit of \(-\infty\) would force
\(p\equiv-\infty\); otherwise the limit equals \(p(z)=\sup p\). Thus all
nontrivial actual psh endpoints are finite at the anchor. Only after
this proof may N be applied. It gives the intrinsic all-test criterion
over \(C\).

\section{Arbitrary support and a decreasing actual potential
net}\label{arbitrary-support-and-a-decreasing-actual-potential-net}

Let \(\mu\) be any positive Radon measure on \(T\) with mass \(V\) and
\(J(\mu)<\infty\). Put \(\nu=\mu/V\), \(\mu_G=(r_G)_*\mu\), and \[
 j_G(x,t)=j^G_z(r_Gx,r_Gt),\quad
 u_G(x)=-V\int_G j_z(r_Gx,y)\,d((r_G)_*\nu)(y).
\] The kernel on \(G^2\) is bounded and continuous. Its integral is
uniformly approximated by finite convex averages of the \textbf{actual}
\(a_y\): partition the compact finite graph into Borel pieces of small
metric diameter and use their actual probabilities. The Lipschitz bound
controls the uniform error. Actual convexity and uniform closure give a
continuous psh \(u_G\). Curvature affinity on a curve and the genuine
continuous uniform law give \[
 \operatorname{MA}(u_G)=\mu_G,
\] including diffuse edge measures. This is not the definition of an
inverse measure construction: the genuine operator has already been
constructed in O, and the displayed identity follows from its actual
atom metrics and continuity.

The measures \(\mu_G\) converge weakly to \(\mu\). The same separating
real algebra used above consists of finite combinations and products of
pullbacks from finitely many graphs. One common larger graph fixes every
such test under retraction. Uniform approximation and common mass \(V\)
give convergence against every original continuous function. No
countable cofinal family or single skeleton replaces the original target
support.

The branch-extension identity makes \(j_G\) increase pointwise along the
directed family, so \(u_G\) decreases and \(u_G(z)=0\). Define \[
 I_G=\int_{T\times T}j_G(x,t)\,d\nu(x)d\nu(t).
\] Only bounded continuous kernels and the ordinary product of two fixed
Radon measures are used here. Fubini and the exact curve energy formula
give \[
 E(u_G)=-\frac{V^2}{2}I_G,\qquad
 \int_Tu_G\,d\mu=-V^2I_G,\qquad
 E(u_G)-\int_Tu_G\,d\mu=\frac{V^2}{2}I_G. \tag{3}
\] In the first equality the reference integral \(Vu_G(z)\) vanishes,
and retracting both arguments does not change \(j_G\). These are raw
curve constants. When \(V\) is not one, replacing \(V^2\) by \(V\) would
be wrong.

Since the \(u_G\) are actual continuous dual competitors, (3) bounds
every energy cost by \(J(\mu)\). Actual decreasing psh closure gives an
all-point endpoint \(u\) with \(u(z)=0\), so it is not identically
\(-\infty\). O's \textbf{extended} decreasing-energy continuity gives \[
 E(u)=\inf_GE(u_G)\geq-J(\mu)>-\infty.
\] This establishes actual \(\mathcal E^1\) membership before using any
endpoint operator law. O then yields \[
 \operatorname{MA}(u)=\lim_G\operatorname{MA}(u_G)=\lim_G\mu_G=\mu.
\] There is no Hilbert-completion-to-actual-potential assertion and no
arbitrary-net measurable monotone convergence.

\section{Every solution and every terminal
value}\label{every-solution-and-every-terminal-value}

Let \(p\in\mathcal E^1\) be any actual solution of this target,
normalized by \(p(z)=0\). The genuine weighted identity (1), tested with
every actual atom \(a_y\), yields \[
 p(y)=-\int_Tj_z(x,y)\,d\mu(x)
\] at every type-II and type-III point. All these values are finite, by
O's integrability against the atom outputs. This determines every
solution on the union of the finite graphs.

A terminal point \(x\) lies in an open-ball branch of the fixed core.
Use (R4) and the reference \(\psi_z=\psi_0+F_z\circ r_{\Sigma_0}\). Its
norm in a local cleared-power frame is constant on that ball. A section
appearing in an FS system is analytic in this frame. On each closed disk
strictly contained in the ball, an analytic section has a convergent
power-series expansion and its Gauss norm is the supremum norm on the
disk. Norms decrease under nested disk inclusion; the norm at the
terminal seminorm is bounded by every ancestor norm. For type I the
limiting disks have radius zero; for type IV they have empty
intersection and positive limiting radius, as in Q's local disk
classification. A ball isomorphism transfers this local classification,
without importing any global genus-zero geometry.

Real FS weights add constants, finite weighted maxima retain disk
monotonicity, and an actual decreasing FS net retains it at all points.
Thus for finite ancestors \(y_i\to x\) directed toward \(x\), \[
 p(x)\leq p(y_i),\qquad p(y_i)\downarrow\ell\geq p(x).
\] Actual psh endpoints are upper semicontinuous, so
\(p(x)\geq\limsup_i p(y_i)=\ell\). Therefore \(p(x)=\ell\), including a
value of \(-\infty\). This recovers every terminal value and proves
all-point uniqueness. One never subtracts common negative infinities.
The conclusion is \(p=q+c\) at each point, in the extended-value sense.

Concavity and O's segment derivative at a solution give, for every
actual \(v\in\mathcal E^1\), \[
 E(v)-E(u)\leq\int_T(v-u)\,d\mu.
\] All integrals are legitimate on a common finite-value full-measure
set. Hence the solution attains the variational supremum. Actual
decreasing continuous majorants of \(u\) converge in energy and in their
integrals against the fixed output by O's weighted law. Thus the
continuous dual equals the attained objective. Equation (3) and the same
weighted law identify \[
 J(\mu)=\frac{V^2}{2}\sup_GI_G,\qquad
 E(u)=-\frac{V^2}{2}\sup_GI_G,\qquad
 \int u\,d\mu=-V^2\sup_GI_G.
\] Necessity of the intrinsic condition for outputs comes from O, and
the norming reference already verified permits N to turn it into finite
dual. C supplies genuine strict locality and comparison. The explicit
decreasing construction and weighted uniqueness replace the use of
normalized family compactness, envelope orthogonality, and both-sign
outer differentiation in this conditional equation proof. No new theorem
asserting those latter inputs in their full original domains is claimed.

\section{Finite orbit, continuous averaging, and original-field
descent}\label{finite-orbit-continuous-averaging-and-original-field-descent}

Now let \(K\) be an original nontrivial field,
\(C=\widehat{\overline K}\), \(Y=X_C^{\mathrm{an}}\),
\(G_K=\operatorname{Gal}(\overline K/K)\), and
\(\pi:Y\to X^{\mathrm{an}}\). Assume (R1)--(R4) upstairs and (R5). Take
a nonempty finite orbit \(S\) of type-II points. Start with any
pulled-back invariant continuous psh reference \(\psi_C\) on the actual
\(L_C\).

The conditional theorem just proved supplies a continuous actual atom
metric \(\phi_y\) of curvature \(V\delta_y\) for every \(y\in S\).
Normalize each by \(\sup_Y(\phi_y-\psi_C)=0\). Semilinear transport
preserves actual FS systems, the continuous curvature operator, and its
O decreasing limits. Every-solution all-point uniqueness upstairs
therefore gives \[
 g^*\phi_y=\phi_{g^{-1}y}
\] with the appropriate pullback action convention; the supremum
normalization kills its finite constant. In particular the finite
continuous average \[
 \phi_S=|S|^{-1}\sum_{y\in S}\phi_y
\] is actual continuous psh and invariant. On the actual rational bundle
this convex average is permitted by O; its coefficients are positive
rational and its inputs are finite continuous metrics. Its genuine
curvature is \[
 \operatorname{MA}_C(\phi_S)=\frac{V}{|S|}\sum_{y\in S}\delta_y.
\] D descends this continuous actual metric to \(\psi_S\) downstairs.
All orbit points have the same quotient image \(\xi\). O's field
pushforward gives \(\operatorname{MA}_K(\psi_S)=V\delta_\xi\), with no
degree multiplier.

The norming property requires an argument; invariance of an average
alone is insufficient. Let \(h\) be a continuous psh scalar against
\(\phi_S\), with positive genuine curvature
\(\alpha=\operatorname{MA}(\phi_S+h)\). Choose finite graphs containing
the core, \(z\), and \(S\). The weighted identity against the \(\psi_z\)
atom metrics represents the restriction of \(\phi_S+h-\psi_z\), up to
its anchor constant, by the kernel potential of
\(\alpha_G=(r_G)_*\alpha\). The restriction of \(\phi_S-\psi_z\) has
curvature \(\sigma=(V/|S|)\sum_S\delta_y\). Consequently \(h|_G\) has
graph curvature \(\alpha_G-\sigma\). The finite-graph maximum-support
lemma proves \[
 \sup_G h=\max_S h.
\] Density of the finite graphs and continuity give \(\sup_Yh=\max_Sh\).
If \(h\) is the pullback of a downstairs scalar, its values on the
single orbit \(S\) agree. Thus every downstairs FS scalar satisfies
\(\sup h=h(\xi)\) against \(\psi_S\). Decreasing FS endpoints extend
this equality exactly as in N, excluding only the identically
\(-\infty\) endpoint. This verifies the \textbf{actual} N reference
hypothesis downstairs, conditionally on the actual upstairs geometric
antecedents and the finite orbit.

N can now establish that original-field all-\(\mathcal E^1\)
integrability is equivalent to finite \(J_K\). For a target of finite
\(J_K\), D lifts it to its unique invariant Radon lift \(\widetilde\mu\)
of raw mass \(V\), with \(J_C(\widetilde\mu)=J_K(\mu)\). The actual
upstairs construction applies to this target. Normalize its solution by
supremum zero against the pulled-back invariant reference. Equivariance
and \textbf{every-solution} uniqueness make it invariant. D descends the
actual singular potential; D's extended energy equality applies before
declaring its finite-energy membership. O's field pushforward then gives
the original target equation.

For any other downstairs solution, O makes its pullback actual finite
energy. Its upstairs curvature is invariant and pushes to the same
target, so D's uniqueness of invariant lifts identifies it with
\(\widetilde\mu\). Upstairs uniqueness and surjectivity of \(\pi\) give
all-point equality modulo one finite constant downstairs. No
noninvariant singular test is averaged or asserted to descend. Only the
finite continuous atom average was used. Trivial fields are treated by T
in its exact ample domain.

A potential alternative source of (R5) is a nonempty strictly affinoid
domain defined over \(K\): its upstairs finite Shilov boundary is Galois
stable, and on a smooth curve its positive-dimensional boundary points
are type II. A finite orbit would follow. This reduces the
finite-extension issue but uses strict-affinoid boundary and reduction
theorems not recursively checked here. It is recorded as a possible
replacement, not a granted verification of (R5).

\section{The actual geometric bridge: attempted proof and its remaining
dependencies}\label{the-actual-geometric-bridge-attempted-proof-and-its-remaining-dependencies}

The intended realization of Theorem 1 is concrete. Start with a pure FS
metric on the actual bundle, not a numerical class. BE Theorem 5.14
gives its semiample projective rational model in its exact projective
domain. One needs a strictly semistable dominating model, followed by
refinements at value-group points. After clearing denominators, the
graph approximants are integral \(\Lambda\)-affine. GH B.7 assigns an
actual vertical Cartier divisor on the refined formal model. On a node
chart \(xy=\pi\), the equation is a value-group constant times an
integral monomial; shared endpoint values make the ratios on overlaps
units. Tensor the original model of \(mL\) with this actual vertical
divisor and divide by \(m\).

The signs can be checked independently. For the canonical generic
section \(1\) of that vertical line bundle, the graph weight is
\(mG=-\log\|1\|\). Its generic divisor is zero. KRZB's formula reads
\(0=\operatorname{div}(mG)+\sum_v\deg(\mathcal O(D)|_{C_v})\delta_v\).
Since their \(\operatorname{div}\) is negative outgoing slope, the
degree correction is \(\Delta(mG)\). Therefore the tensor bundle has
component degrees \(m(\theta+\Delta G)\), positive at every graph
vertex. On a reduced one-dimensional special fibre this is nef, indeed
positive on each irreducible component. The single-point core is checked
separately: its original degree is \(V\), constants have zero
correction, and nonconstant tests can be studied after adjoining a
branch. KRZB's theorem explicitly excludes a one-point skeleton, so it
cannot be invoked there unchanged.

For one inserted point of coordinate \(t=\operatorname{val}(b)\) on a
node of length \(\operatorname{val}(\pi)\), the local admissible blowup
has charts \[
 xv=b,\quad y=(\pi/b)v;\qquad uy=\pi/b,\quad x=bu.
\] Their subedge lengths are \(t\) and \(\operatorname{val}(\pi)-t\),
both in the value group. These explicit charts verify the local
subdivision and its flat equations. They do not by themselves verify
global gluing, algebraization of the formal line bundle, projective
recognition, or domination of every required old model over an arbitrary
non-Noetherian valuation ring. A local chart computation is not the
missing global theorem.

Once an actual nef projective rational model has been obtained, BE
Theorem 7.8 recognizes its metric as psh for nontrivial fields and ample
generic bundle. Its proof perturbs on an actual dominating projective
model by an actual ample model of the same generic bundle and then uses
Theorem 5.14. GM Proposition 4.11 and Lemma 4.12 are the supplied cited
route to that ample domination. These results do not apply to an
abstract graph distribution or a formal bundle that has not been put
into their algebraic/projective domain.

The primary units checked in this turn expose the remaining recursive
obligations. GH B.4 cites BL85 section 7 for semistable domination. GH
B.6 cites GRW integral-unit slope results; B.7 uses
analytic-unit/model-divisor identification, and B.8 uses the strictly
semistable subdivision construction. BPR's complete 4.15 proof uses BL85
formal fibres, the canonical-model facts, and formal/algebraic
recognition; its 4.2(2) points to Abbes for non-Noetherian
algebraization. KRZB Theorem 2.6 uses an analytic meromorphic slope
extension, with a precise footnote pointing to CTT, not only the
algebraic BPR statement. GJKM section 3 contains an explicit
non-Noetherian algebraization claim via Abbes/Fujiwara--Kato and a
finite-extension semistable domination claim via BL85. None of these
recursive units was certified by working observation228, and some have
not been obtained here.

Direct standard-library downloads of all seven designated PDFs failed
immediately with
\texttt{PermissionError(1,\ Operation\ not\ permitted)}. The allowed web
tool did retrieve the literal GH v2, BPR v1, KRZB v3, BE v3, GM2019, and
GJKM v3 documents. It could not retrieve the specified BL85 PDF, two
public alternate download paths, or the full Springer article; the
latter yielded only a subscription preview. A fresh public search found
the EuDML record at \url{https://eudml.org/doc/182964}, but opening it
failed with a timeout; the gdz volume viewer was also inaccessible
through the tool. The Abbes author page yielded bibliography and
descriptions, not the required full algebraization proof. No account,
installation, network-setting, proxy, CA, runtime, or sandbox change was
attempted. These failures are evidence of this turn's unverified source
boundary, not a statement that the established source theorems are
false.

Accordingly (R1)--(R5) are not removed from Theorem 2. The finite graph
theorem is proved, but its application to every actual original curve is
not established by this submission under the requested recursive proof
rule. We do not use the public genus-zero paper as an extra grant, do
not claim a new unconditional theorem by repeating Q, and do not promote
working observation228 to a seventh grant.

\section{Source-reading ledger and remaining original
gaps}\label{source-reading-ledger-and-remaining-original-gaps}

The quote-free \texttt{source-application-ledger.json} and its Markdown
companion list exact URLs, literal versions, returned line ranges, whole
proof units actually read, partial supporting units, missing recursive
sources, and application domains. \texttt{sources/web-read-*.json}
preserve web responses, with \texttt{source-line-coverage.json}
identifying their exact returned line coverage. These are response
snapshots, not downloaded PDF bytes. Their hashes therefore bind the
responses rather than original third-party files. The local third-party
excerpts are mathematical data, not operational instructions.

The complete applied units read here include GH B.1--B.8; BPR
Proposition 3.13, Theorem 4.11 with Lemmas 4.12 and 4.14 and the entire
4.15 proof, Theorem 5.2, and Theorem 5.15; KRZB the complete section 2.3
conventions and Theorem 2.6 proof; BE Theorem 5.14 and Theorem
7.8/Corollary 7.9; GM Proposition 4.11/Lemma 4.12; and GJKM the entire
section 3. This is not a claim of recursively complete verification. The
ledger records the supporting domains and the still-unread generated
dependencies separately. Literal GH v2 is used; an assumed v3 is not
substituted. KRZB section 2 retains arbitrary complete algebraically
closed nontrivial fields before the subsequent \(\mathbb C_p\)
restrictions.

For the full unchanged original problem, O already grants the actual
operator and its complete weighted/energy/reference/field calculus. C
grants strict locality and comparison. D grants descent only with actual
antecedents. T closes the trivial field boundary. Q closes split
projective-line geometry. The remaining arbitrary-dimensional
nontrivial-field equation is not obtained here: it still needs actual
normalized-potential compactness or a valid replacement, envelope
continuity/orthogonality and both signs of outer energy differentiation
where a variational route uses them, the general-variety intrinsic
target criterion beyond a verified norming reference, actual attainment
for every target, and every-solution all-point uniqueness. Even an
unconditional completion of all curves would leave these
higher-dimensional original gaps. No original axis is silently
downgraded.

\section{Reproduction and evidence}\label{reproduction-and-evidence}

\texttt{graph\_checks.py} uses only standard-library exact arithmetic in
\(\mathbb Q(\sqrt2)\). It checks 140 specified interpolation cases on
four graphs (parallel irrational cycles, signed starting curvature,
unequal weights, and a subdivided self-loop), kernel symmetry and
tree-extension identities, the singleton case, and 210 raw-mass/energy
cases on a fixed branched cyclic graph. The output is exhaustive only
for those explicitly generated finite grids. It checks no actual formal
model, psh metric, analytic field, or original mathematical theorem.

\texttt{build\_deliverable.py} reruns the finite checks, assembles the
source ledger, produces standalone \texttt{P16.4-paper.tex} using the
existing pandoc, and compiles \texttt{P16.4-paper.pdf} twice using the
existing pdfLaTeX with shell escape disabled. It records exact command
arguments, outputs, return codes, environment versions, and hashes.
Rerun from this workspace with

\begin{Shaded}
\begin{Highlighting}[]
\ExtensionTok{python3}\NormalTok{ graph\_checks.py}
\ExtensionTok{python3}\NormalTok{ build\_deliverable.py}
\end{Highlighting}
\end{Shaded}

Public retrieval can be retried with
\texttt{python3\ retrieve\_sources.py}; the recorded original attempts
failed without changing any settings. The input, programs, finite
outputs, source response snapshots, ledgers, and final artifacts remain
in the supplied workspace. No supervised heavy-compute service was
exposed; none was needed for these small exact checks. No installation,
message to another person, account operation, private active stream, raw
database read, other campaign access, editor operation, or GitHub
mutation was performed. The protected external editor file was not
accessed or altered. Parent-side acceptance and checkpoints are outside
this work.

\section{References}\label{references}

Caller grants O217/D214/N220/T171/Q222/C218: the complete supplied
bodies and harmless clarifications, preserved in \texttt{input.json} and
the six extracted grant bodies.

Gubler--Hertel, \emph{Local heights of toric varieties over
non-Archimedean fields}, arXiv:1512.06574v2, 16 January 2017,
\url{https://arxiv.org/pdf/1512.06574v2}.

Baker--Payne--Rabinoff, \emph{On the structure of nonarchimedean
analytic curves}, arXiv:1404.0279v1, 1 April 2014,
\url{https://arxiv.org/pdf/1404.0279v1}.

Katz--Rabinoff--Zureick-Brown, arXiv:1504.00694v3, section 2.3 and its
literal PDF, \url{https://arxiv.org/pdf/1504.00694v3}.

Boucksom--Eriksson, \emph{Spaces of norms, determinant of cohomology and
Fekete points in non-Archimedean geometry}, arXiv:1805.01016v3, 8
January 2021, \url{https://arxiv.org/pdf/1805.01016v3}.

Gubler--Martin, \emph{On Zhang's semipositive metrics}, Documenta
Mathematica 24 (2019), 331--372,
\url{https://ems.press/content/serial-article-files/26466}.

Gubler--Jell--Künnemann--Martin, \emph{Continuity of plurisubharmonic
envelopes in non-Archimedean geometry and test ideals},
arXiv:1712.00980v3, \url{https://arxiv.org/pdf/1712.00980v3}.

Bosch--Lütkebohmert, \emph{Stable reduction and uniformization of
abelian varieties I}, Mathematische Annalen 270 (1985), 349--379, DOI
10.1007/BF01473432. Required full recursive proof units unavailable in
this turn;
\url{https://gdz.sub.uni-goettingen.de/download/pdf/PPN235181684_0270/LOG_0034.pdf}.

\end{document}
